\documentclass[11pt]{article}
\usepackage[margin=1in]{geometry}
\usepackage{amsmath,amssymb,amsthm}
\usepackage{booktabs,tabularx,array}
\usepackage[hidelinks]{hyperref}

\newtheorem{definition}{Definition}
\newtheorem{proposition}{Proposition}
\newtheorem{corollary}{Corollary}
\newcommand{\NT}{\mathsf{NT}}
\newcommand{\Ready}{\mathsf{Ready}}
\newcommand{\Review}{\mathsf{Review}}
\newcommand{\Promote}{\mathsf{Promote}}
\newcommand{\Pos}{\mathsf{Pos}}

\title{P14-T03 No-Target and Positive-Provenance Sufficiency Policy v1}
\author{The AEther-Flow Research-Control Project}
\date{July 22, 2026}

\begin{document}
\maketitle

\begin{abstract}
This task-local draft/control policy fixes the logical role of a no-target
certificate.  Source purity is necessary for a physical-readiness record, but
it is not positive provenance and cannot by itself establish physical meaning,
emergence, Gate A through D readiness, or promotion.  Readiness additionally
requires gate-scoped evidence for source derivation, uniqueness or an explicit
quotient, naturality, dynamics, operational systems, robustness, and
appropriately provenanced independent review.  Even an evidence-complete record
requires a separately authorized protected verdict.  The policy preserves the
current status of all inspected objects and does not reject source-pure
mathematics.
\end{abstract}

\section{Authority and source boundary}

This artifact executes only v21 P14-T03 under task RT-20260722-019.  It is not
canonical ontology, a source law, a physical gauge theorem, a Gate Chair
decision, proof authority, a scientific-ledger change, or a completed
first-principles derivation of general relativity.  Its source basis is the
registered v21 implementation plan, the P11-T05 positive-provenance policy, the
P14-T01 success-category theorem, and the registered foundations and geometry
manuscripts.  No external source is used.

The foundations and geometry manuscripts distinguish exact-GR \emph{adoption}
from substrate \emph{derivation}.  P14-T01 separately types interpretive
redescription, formal or categorical equivalence, genuine emergence, and
empirical novelty.  P11-T05 supplies seven positive-evidence dimensions and
states that neither a no-target certificate nor an operational validator
receipt completes a physical-readiness record.

\section{Logical policy}

Let $x$ be a typed, scoped claim record and let $g\in\{A,B,C,D\}$ be a
declared readiness gate.  Every predicate below is evaluated only in the scope
and against the source hashes carried by $x$.

\begin{definition}[No-target certificate]
$\NT(x)$ means that an audit has shown, within its declared search and
variation scope, that the construction does not use a prohibited target
benchmark, target atlas, target metric, validator result, checkpoint, or test
receipt as a source premise.  It certifies structural admissibility in that
scope.  It neither supplies a source derivation nor assigns physical meaning.
\end{definition}

\begin{definition}[Positive-evidence vector]
The positive-evidence vector is
\[
 \Pos_g(x)=(D_s,U_q,N,D_y,O,R,I)_g,
\]
where the coordinates respectively denote gate-adequate certificates for
source derivation, uniqueness or a defined quotient, naturality, dynamics,
operational systems, robustness, and independent review.  A coordinate is true
only when its typed evidence meets the minimum status for gate $g$; missing,
not-applicable, assumed, proposed, conditional, or merely specified evidence
does not silently become complete.
\end{definition}

\begin{definition}[Evidence completeness and protected promotion]
Let $S_g(x)$ mean that the record is well typed for gate $g$, has exact source
and scope identity, and satisfies the gate-specific burdens listed below.  Set
\[
 \Ready_g(x) := \NT(x)\wedge S_g(x)\wedge
 D_s(x)\wedge U_q(x)\wedge N(x)\wedge D_y(x)\wedge
 O(x)\wedge R(x)\wedge I(x).
\]
$\Ready_g(x)$ means only ``evidence complete for a separately authorized gate
review.''  If $V_g(x)$ denotes a valid protected verdict with exact authority,
then
\[
 \Promote_g(x) := \Ready_g(x)\wedge V_g(x).
\]
Operational validators may check the shape and internal consistency of these
records, but do not make any conjunct true as scientific evidence.
\end{definition}

\begin{proposition}[NT-POSITIVE-SUFFICIENCY-THEOREM-V1]
For every well-typed record $x$ and gate $g$:
\begin{enumerate}
  \item $\Ready_g(x)\Rightarrow\NT(x)$;
  \item $\NT(x)\not\Rightarrow\Ready_g(x)$;
  \item $\Ready_g(x)\not\Rightarrow\Promote_g(x)$; and
  \item no individual positive coordinate is supplied by $\NT(x)$.
\end{enumerate}
Thus no-target purity is necessary but insufficient for evidence completeness,
and evidence completeness is itself insufficient for protected promotion.
\end{proposition}

\begin{proof}
The first implication is immediate from the conjunction defining
$\Ready_g$.  For the second, take the finite evidence world
\[
 w_{\NT}=(\NT,S_g,D_s,U_q,N,D_y,O,R,I,V_g)
        =(1,1,0,0,0,0,0,0,0,0).
\]
It has a valid no-target certificate but no positive-evidence certificate, so
$\Ready_g=0$.  For the third, take
\[
 w_{\Review}=(1,1,1,1,1,1,1,1,1,0).
\]
This record is review-ready, but $V_g=0$, hence $\Promote_g=0$.  Finally, for
each positive coordinate $P$ use the two worlds $w_{\NT}$ and
$w_{\NT,P}$, the latter differing only by $P=1$.  Both satisfy $\NT$, while
only one satisfies $P$, so $\NT\not\Rightarrow P$.  These are finite logical
countermodels to overreading a purity certificate; they are not models of
nature and not no-go theorems about future source extensions.
\end{proof}

\begin{corollary}[Source-pure mathematics remains admissible]
If $\NT(x)=1$ while one or more positive coordinates are missing, $x$ may
remain valid scoped mathematics, a proposal-only construction, source-extension
data, or a draft/control theorem.  The policy blocks only the stronger physical
or protected conclusion; it does not reject the mathematical result.
\end{corollary}

\section{Gate A through D evidence mapping}

\begin{center}
\small
\begin{tabularx}{\textwidth}{>{\raggedright\arraybackslash}p{0.10\linewidth}>{\raggedright\arraybackslash}p{0.27\linewidth}>{\raggedright\arraybackslash}X>{\raggedright\arraybackslash}p{0.18\linewidth}}
\toprule
Gate & Gate-specific burden & How the positive vector applies & Blocked overread \\
\midrule
A & Source ontology, objects and morphisms, dynamics, EqSrc or quotient status, variations, observables, automorphisms, audit, stress, and review. & Source derivation, uniqueness/quotient, naturality, dynamics, robustness, and review must be evidenced at source scope; operational systems must identify source observables or protocols rather than target labels. & No metric or downstream promotion. \\
B & Derived causal/conformal structure, calibrated scale, signature, nondegeneracy, transformation law, uniqueness up to declared equivalence, robustness, and propagation meaning. & All seven coordinates use the P11-T05 Gate B ready statuses.  A target-independent tensor formula is not yet an effective physical metric. & No matter coupling or field equations. \\
C & Source matter ontology and dynamics, devices, common propagation geometry, universal coupling, action, stress-energy, conservation, equivalence-principle status, audit, stress, and review. & All seven coordinates use the stricter P11-T05 Gate C statuses, including validated dynamics and operational systems. & No Einstein equations or benchmark promotion. \\
D & Selected closure route, satisfied assumptions, derived action or constraints, field equations, modes, hyperbolicity, stability, corrections, audit, and review. & Source provenance must extend through geometry and matter inputs; uniqueness/quotient, dynamics, operational meaning, robustness, and review remain explicit rather than inherited from exact-GR agreement. & No benchmark or completed-derivation promotion. \\
\bottomrule
\end{tabularx}
\end{center}

No Gate A through D verdict is issued here.  Each gate retains its own protected
authority and a narrow result at one gate does not satisfy a downstream gate.

\section{Status-preserving applications}

\begin{center}
\small
\begin{tabularx}{\textwidth}{>{\raggedright\arraybackslash}p{0.20\linewidth}>{\raggedright\arraybackslash}p{0.24\linewidth}>{\raggedright\arraybackslash}X}
\toprule
Object class & Preserved status & Policy consequence \\
\midrule
$M_{\rm src}$ & Scoped source-only adopted object. & Its name or source scope does not by itself establish a completed source theory, a unique physical realization, dynamics, operational meaning, or Gate A readiness. \\
Scoped metric records & Scoped source-extension objects, not an unscoped Lorentzian $g_{\rm eff}$. & A no-target or source-purity result does not establish Gate B's causal, scale, signature, uniqueness, operational, robustness, and review burdens. \\
Matter profiles & Scoped evidence and preconditions, not physical matter coupling. & Source-pure profiles do not construct devices, universal coupling, an action, stress-energy, conservation, or Gate C readiness. \\
Selector and quotient candidates & Proposal-only or draft/control in their exact recorded scopes; scoped negative results remain scoped. & Equivariance or target independence does not establish a physical gauge, source-law adoption, unique selection, operational semantics, or general EqSrc completion. \\
Exact-GR benchmark line & Interpretive redescription in the P14-T01 mapping. & Exact recovery remains target-side benchmark evidence; it is not source derivation, emergence, empirical novelty, or Gate D closure. \\
\bottomrule
\end{tabularx}
\end{center}

The machine-readable examples accompanying this paper do not change any
canonical status.  ``Not assessed'' remains distinct from ``failed,'' and an
unmet present evidence obligation is not an impossibility theorem.

\section{Overread controls and theorem-inventory guidance}

The following formulations are blocked unless their independent evidence is
present: ``target independent, therefore physical''; ``source pure, therefore
emergent''; ``equivariant, therefore gauge''; ``validator passed, therefore
proved''; ``exact-GR, therefore derived''; and ``evidence complete, therefore
adopted.''

Future theorem-inventory rows should record no-target status separately from
the seven positive coordinates, the gate and scope, exact source hashes,
assumption and quotient boundaries, operational protocol references,
robustness domain, review provenance, scientific status, and protected
authority.  A no-target PASS must never auto-fill a positive coordinate.  A
finite countermodel or scoped obstruction must retain its scope and reopening
conditions rather than becoming a global theory rejection.

\section{Conclusion}

The policy supplies a claim-logic theorem and finite countermodels, not a new physical law.
It closes the P14-T03 policy burden at draft/control scope while
leaving all physical derivation routes open.  Canonical ontology, scientific
status, Distance-to-GR ledgers, Gate Chair authority, proof authority,
publication authority, and benchmark or completed-derivation status remain
unchanged.

\end{document}
